\documentclass[12pt]{article}
\usepackage{amsmath,amssymb,amsthm}
\usepackage[margin=1in]{geometry}

\newtheorem{theorem}{Theorem}
\newtheorem{lemma}[theorem]{Lemma}

\title{Project Euler Problem 501: Eight Divisors}
\author{}
\date{}

\begin{document}
\maketitle

\section{Problem Statement}
Let $f(n)$ be the count of numbers not exceeding $n$ with exactly eight divisors.
Given: $f(100) = 10$, $f(1000) = 180$, $f(10^6) = 224427$. Find $f(10^{12})$.

\section{Mathematical Foundation}

\begin{theorem}[Divisor Count Formula]
Let $n = p_1^{a_1} p_2^{a_2} \cdots p_k^{a_k}$. Then $\tau(n) = \prod_{i=1}^{k}(a_i + 1)$.
\end{theorem}

\begin{proof}
Each divisor of $n$ has the form $\prod p_i^{b_i}$ with $0 \le b_i \le a_i$. There are $a_i + 1$ independent choices for each $b_i$, so by the multiplication principle, $\tau(n) = \prod(a_i + 1)$.
\end{proof}

\begin{lemma}[Classification of $\tau(n) = 8$]
A positive integer $n$ satisfies $\tau(n) = 8$ if and only if $n$ takes one of the forms:
\begin{enumerate}
    \item $n = p^7$ for some prime~$p$,
    \item $n = p^3 q$ for distinct primes $p, q$,
    \item $n = pqr$ for distinct primes $p < q < r$.
\end{enumerate}
\end{lemma}

\begin{proof}
We need $\prod(a_i + 1) = 8 = 2^3$. The multiplicative partitions of~8 into factors $\ge 2$ are $\{8\}$, $\{4, 2\}$, and $\{2, 2, 2\}$, yielding the three forms above. These are clearly mutually exclusive and exhaust all possibilities.
\end{proof}

\begin{theorem}[Counting Formula]
Let $N = 10^{12}$ and $\pi(x)$ denote the prime counting function. Then
\[
f(N) = C_1 + C_2 + C_3
\]
where:
\begin{align}
C_1 &= \pi(N^{1/7}), \\
C_2 &= \sum_{\substack{p \text{ prime} \\ p^3 < N}} \left[\pi\!\left(\frac{N}{p^3}\right) - \mathbf{1}[p^4 \le N]\right], \\
C_3 &= \sum_{p \text{ prime}} \sum_{\substack{q \text{ prime} \\ p < q,\; pq^2 \le N}} \left[\pi\!\left(\frac{N}{pq}\right) - \pi(q)\right].
\end{align}
\end{theorem}

\begin{proof}
\textbf{Case 1:} $n = p^7 \le N$ iff $p \le N^{1/7}$, giving count $\pi(N^{1/7})$.

\textbf{Case 2:} $n = p^3 q \le N$ with $p \ne q$. For each prime~$p$ with $p^3 < N$, the number of valid primes $q \le N/p^3$ is $\pi(N/p^3)$, minus~1 when $q = p$ is itself valid (i.e., when $p^4 \le N$).

\textbf{Case 3:} $n = pqr$ with $p < q < r$. For each pair $(p,q)$ with $p < q$, valid primes $r$ satisfy $q < r \le N/(pq)$, giving $\pi(N/(pq)) - \pi(q)$ choices. We need $pq^2 < N$ for at least one valid~$r$.

By Lemma~1 the three cases are disjoint and exhaustive.
\end{proof}

\section{Algorithm}
\begin{verbatim}
1. Sieve primes up to sqrt(N) = 10^6.
2. Implement pi(x) via Lucy Hedgehog / Meissel-Lehmer.
3. Case 1: C1 = pi(floor(N^(1/7))).
4. Case 2: For each prime p with p^3 < N:
     C2 += pi(N/p^3) - [p^4 <= N].
5. Case 3: For each prime p, for each prime q > p
     with p*q^2 <= N:
     C3 += pi(N/(p*q)) - pi(q).
6. Return C1 + C2 + C3.
\end{verbatim}

\section{Complexity Analysis}
\begin{itemize}
    \item \textbf{Time:} $O(N^{2/3})$, dominated by the prime counting function.
    \item \textbf{Space:} $O(N^{1/2})$ for the sieve and $\pi(x)$ tables.
\end{itemize}

\section{Answer}

\[
\boxed{197912312715}
\]

\end{document}
